\documentclass{article}
\usepackage{svg}
\usepackage{amsmath}
\usepackage[utf8]{inputenc}
\usepackage[french]{babel}
\usepackage{graphicx}
\usepackage[T1]{fontenc}
\usepackage{hyperref}
\hypersetup{
    colorlinks,
    citecolor=black,
    filecolor=black,
    linkcolor=black,
    urlcolor=black
}
\graphicspath{ressources/}
\usepackage{tikz}
\usetikzlibrary{shapes,positioning}

\usepackage{blindtext}

\usepackage{subfiles} % Best loaded last in the preamble

\title{Titre}
\author{ }
\date{ }

\begin{document}
\begin{titlepage}
    \begin{center}
        
        {\Large Université de Mons}\\[1ex]
        {\Large Faculté Des Sciences}\\[1ex]
        {\Large Département d'Informatique}\\[1ex]
        
        \newcommand{\HRule}{\rule{\linewidth}{0.3mm}}
        % Title
        \HRule \\[0.3cm]
        { \LARGE \bfseries Projet de modélisation logicielle \\[0.3cm]}
        { \LARGE \bfseries Rapport de modélisation  \\[0.1cm]} % Commenter si pas besoin
        \HRule \\[1.5cm]
        
        % Author and supervisor
        \begin{minipage}[t]{0.45\textwidth}
            \begin{flushleft} \large
                \emph{Professeur:}\\
                Tom \textsc{Mens}\\
                \emph{Assistants:}\\
                Jeremy \textsc{Dubrulle}\\
                Gauvin \textsc{Devillez}\\
                Sébastien \textsc{Bonte}\\
            \end{flushleft}
        \end{minipage}
        \begin{minipage}[t]{0.45\textwidth}
            \begin{flushright} \large
                \emph{Auteurs:} \\
                Pignozi \textsc{Agbenda} \\
                Thomas \textsc{Bernard} \\
                Théo \textsc{Godin} \\
                Ugo \textsc{Proietti}\\
            \end{flushright}
        \end{minipage}\\[2ex]
        
        \vfill
        
        % Bottom of the page
        \begin{center}
            \begin{tabular}[t]{c c c}
                \includegraphics[height=1.5cm]{ressources/logoumons.jpg} &
                \includegraphics[height=1.5cm]{ressources/logofs.jpg} &
            \end{tabular}
        \end{center}~\\
        
        {\large Année académique 2021-2022}
        
    \end{center}
\end{titlepage}

\title{Projet de Modélisation Logicielle}
\author{AGBENDA Christian, BERNARD Thomas, GODIN Theo et PROIETTI Ugo}
\maketitle
\tableofcontents
\newpage

\pagenumbering{arabic}

\section{Introduction}

\section{Applications de bases}

	\subsection{Application1 : Gestion de portefeuilles financiers}

		\subsubsection{Vue d'ensemble}
			L'application 1 est l'application qui est destinée aux clients. Il s'agit de l'application sur
			laquelle nous avons commencé à travailler et donc celle qui a subi le plus de modifications au total.
			Il est possible que certains concepts soient modifiés lors de l'implémentation notamment ceux qui 
			n'ont pas été approfondis au maximum. Cette application sera étendue par l'entièreté du groupe.

		\subsubsection{Diagramme de cas d'utilisation}
				
				\subfile{sections/usecase1}
		\newpage
		\subsubsection{Interaction Overview Diagram}
		
				\subfile{sections/rapportInteractionOverviewApp1}
		
		\subsubsection{Diagramme de classes}

		\subsubsection{Diagrammes de séquences}
		
	\subsection{Application 2 : Application de gestion pour insitutions financières}

		\subsubsection{Vue d'ensemble}
			L'application 2 est l'application destinée aux institutions. Cette application leur permet de gérer 
			leurs clients ainsi que les produits de ceux-ci. La conception de cette application est basée sur celle 
			de l'application 1. En effet l'application 2 reprend certains des aspects de l'applications 1. Tout 
			en étant plus restrictive. En effet, l'application 2 n'a pas d'accès à la notion de portfeuille. 
			Toutefois, elle possède des outils de gestion qui ne font pas partie de l'applciation 1.
		\subsubsection{Diagramme de cas d'utilisation}
			\subfile{sections/usecase2}
		\newpage
		\subsubsection{Interaction Overview Diagram}

				\subfile{sections/rapportInteractionOverviewApp2}

		\subsubsection{Diagramme de classes}

		\subsubsection{Diagramme de séquences}
	
	\subsection{Serveur}
	
		\subsubsection{Vue d'ensemble}
		
		\subsubsection{Diagramme d'entité-relation}
			
			\subfile{sections/erd}
		\subsubsection{API}
				\subfile{sections/API}

    \subsection{Interface Graphique}
				\subfile{sections/rapportGui}		

\section{Extensions}
	
	\subsection{Extension 1 : }
	
	\subsection{Extension 2 : }
	
	\subsection{Extension 3 : }
	
	\subsection{Extension 4 : }
	
	
		

\end{document}
